\section{Uniform compression from partitions}
\label{sec:compression}

We first prove an approximation result for an arbitrary family of partitions
of a finite set.  Its cost depends on the number of labels of a partition,
whereas the underlying set can be much larger.

Let $X$ be a finite nonempty set, let $u,q$ be positive integers, and let
\[
  L_i:X\longrightarrow\Sigma_i,\qquad i\in[u]=\{1,\ldots,u\},
  \qquad |L_i(X)|\le q.
\]
Fix $0<\theta<1$ and a nonempty family $\mathcal T$ of subsets of $[u]$
such that
\begin{equation}
  |[u]\setminus T|\le\theta u\qquad(T\in\mathcal T).
  \label{eq:admissible-sets}
\end{equation}
In particular, every $T\in\mathcal T$ is nonempty.  For each such $T$, form
the undirected graph on $X$ in which distinct vertices $x,z$ are adjacent
when $L_i(x)=L_i(z)$ for some $i\in T$.  Write $d_T(x,z)$ for its graph
distance, with distance $\infty$ between different components.

Choose a positive integer $h$ such that
\begin{equation}
  \frac{\log3}{\log(h+1)}\le\theta,
  \qquad s=\lceil1/\theta\rceil.
  \label{eq:compression-parameters}
\end{equation}
Define the logarithmic profile
\[
  \phi_h(d)=\max\left\{0,1-\frac{\log(d+1)}{\log(h+1)}\right\}
  \quad(d=0,1,2,\ldots),\qquad \phi_h(\infty)=0.
\]
Set $Y=X\times\mathcal T$.  For $y=(v,T)\in Y$, define the column function
\begin{equation}
  g_y(x)=\phi_h(d_T(x,v))\qquad(x\in X).
  \label{eq:column-functions}
\end{equation}
Thus $0\le g_y\le1$.  For nonnegative coefficients of total mass at most
one, write
\[
  f_\mu=\sum_{y\in Y}\mu_y g_y,
  \qquad
  \mathcal F_+=\left\{f_\mu:\mu_y\ge0,\ \sum_{y\in Y}\mu_y\le1\right\}.
\]
For a function $f:X\to\R$, let
$\|f\|_{\infty,X}=\max_{x\in X}|f(x)|$.

We seek a uniform approximation list for $\mathcal F_+$ whose size is
controlled by $h,\theta,u,q$, rather than $|X|$.  The next lemma expresses
each profile as a weighted sum of graph-ball indicators and bounds the
error caused by enlarging their radii from $j$ to $3j+2$.

\begin{lemma}[Logarithmic profile]
\label{lem:log-profile}
For any positive integer $h$, the weights
\[
  \gamma_j=\frac{\log(j+2)-\log(j+1)}{\log(h+1)}
  \quad(0\le j<h)
\]
are positive, sum to one, and satisfy, for every
$d\in\{0,1,2,\ldots\}\cup\{\infty\}$,
\begin{align}
  \phi_h(d)
    &=\sum_{j=0}^{h-1}\gamma_j\mathbf1_{\{d\le j\}},
    \label{eq:profile-expansion}\\
  \sum_{j=0}^{h-1}\gamma_j\mathbf1_{\{d\le3j+2\}}
    &\le\phi_h(d)+\frac{\log3}{\log(h+1)}.
    \label{eq:profile-dilation}
\end{align}
\end{lemma}

\begin{proof}
Positivity and the identity $\sum_j\gamma_j=1$ follow by telescoping.
If $0\le d<h$, the sum in \eqref{eq:profile-expansion} runs from $j=d$
to $j=h-1$ and equals
\[
  \frac{\log(h+1)-\log(d+1)}{\log(h+1)}.
\]
If $d\ge h$, including $d=\infty$, both sides are zero.

For finite $d$, the extra indices counted on the left of
\eqref{eq:profile-dilation} satisfy $j<d\le3j+2$.  When this index set is
nonempty, it is the integer interval from
\[
  a_d=\max\left\{0,\left\lceil\frac{d+1}{3}\right\rceil-1\right\}
  \quad\text{to}\quad
  b_d=\min\{h-1,d-1\}.
\]
Since $b_d+2\le d+1$ and $a_d+1\ge(d+1)/3$, its weight is
\[
  \sum_{j=a_d}^{b_d}\gamma_j
  =\frac{\log((b_d+2)/(a_d+1))}{\log(h+1)}
  \le\frac{\log3}{\log(h+1)}.
\]
An empty interval contributes zero.  If $d=\infty$, all the indicators
vanish.  This also covers all distances at and beyond the cutoff $h$.
\end{proof}

\begin{proposition}[Uniform compression]
\label{prop:compression}
Under the hypotheses above, there is a finite list $\mathcal A$ of
real-valued functions on $X$, of cardinality $Q$, such that every
$f\in\mathcal F_+$ has an approximant $A\in\mathcal A$ satisfying
\begin{equation}
  f(x)-3\theta\le A(x)\le f(x)+\theta\qquad(x\in X).
  \label{eq:positive-approximation}
\end{equation}
In particular, $\|f-A\|_{\infty,X}\le3\theta$.  The list can be chosen with
\begin{equation}
  \log Q\le hs\log q+C_{h,\theta,u},
  \qquad
  C_{h,\theta,u}
   =\log u+hs2^u\log\left(1+\frac{s2^u}{\theta}\right).
  \label{eq:compression-count}
\end{equation}
\end{proposition}

\begin{proof}
For each $i\in[u]$, fix a representative map
\[
  \rho_i:L_i(X)\longrightarrow X,
  \qquad L_i(\rho_i(\sigma))=\sigma.
\]
These maps are fixed for the entire approximation problem, before any
coefficient vector is chosen.

\emph{A common index.}
Fix $f=f_\mu\in\mathcal F_+$, and use
$\mu(S)=\sum_{y\in S}\mu_y$ for subsets $S\subseteq Y$.  Averaging
\eqref{eq:admissible-sets} gives
\[
  \frac1u\sum_{i=1}^u\mu\{(v,T):i\notin T\}
   =\sum_{(v,T)\in Y}\mu_{(v,T)}\frac{|[u]\setminus T|}{u}
   \le\theta.
\]
Choose $i_0\in[u]$ for which the omitted mass is at most $\theta$, and
retain the columns in
\[
  Y_0=\{(v,T)\in Y:i_0\in T\}.
\]
This one index will serve at every level $j$.

\emph{A bounded number of pivots at each level.}
Fix $j\in\{0,\ldots,h-1\}$ and, for $x\in X$, let
\[
  I_x=\{(v,T)\in Y_0:d_T(x,v)\le j\}.
\]
Starting with an empty union, add a set $I_z$ whenever its mass outside
the current union is greater than $\theta$.  If $k_j$ sets are added,
their successive contributions give $k_j\theta<1$; in particular
$k_j\le s$.  The process therefore terminates with pivots
$z_{j,1},\ldots,z_{j,k_j}$ and union
\[
  U_j=\bigcup_{k=1}^{k_j}I_{z_{j,k}}
  \quad\text{such that}\quad
  \mu(I_x\setminus U_j)\le\theta\qquad(x\in X).
\]
The empty choice is allowed.  Assign every column of $U_j$ to the first
pivot whose set contains it, and denote its assigned slot by
$\kappa_j(y)$.

For a pivot $z=z_{j,k}$, retain only its label $L_{i_0}(z)$, and replace
it by $z'=\rho_{i_0}(L_{i_0}(z))$.  If $y=(v,T)$ is assigned to this
pivot, then $i_0\in T$, $d_T(v,z)\le j$, and $d_T(z,z')\le1$.
Consequently,
\begin{equation}
  d_T(v,z')\le j+1.
  \label{eq:pivot-replacement}
\end{equation}
For every $x\in X$, the triangle inequality now gives the two implications
\begin{align}
  d_T(x,v)\le j
   &\ \Longrightarrow\ d_T(x,z')\le2j+1,
   \label{eq:pivot-forward}\\
  d_T(x,z')\le2j+1
   &\ \Longrightarrow\ d_T(x,v)\le3j+2.
   \label{eq:pivot-backward}
\end{align}
All distances used in these triangle inequalities are finite under their
respective hypotheses, even if the graph is disconnected.
Figure~\ref{fig:pivot-comparison} illustrates the two bounds.

\begin{figure}[hbp]
\centering
\begin{minipage}[t]{.48\linewidth}
\centering
{\small\textup{(a)} Keeping an assigned match\par}
\medskip
\begin{tikzpicture}[font=\small,
 point/.style={circle,fill=black,inner sep=1.6pt},
 path/.style={draw=black!75,line width=.6pt},
 bound/.style={draw=linkblue,dashed,line width=.7pt}]
\node[point,label=above:$x$] (x) at (0,0) {};
\node[point,label=above:$v$] (v) at (1.5,0) {};
\node[point,label=above:$z$] (z) at (3,0) {};
\node[point,label=above:$z'$] (zp) at (4.5,0) {};
\draw[path] (x) -- node[below] {$\le j$} (v);
\draw[path] (v) -- node[below] {$\le j$} (z);
\draw[path] (z) -- node[below] {$\le1$} (zp);
\draw[bound] (x) to[bend right=40]
 node[below=3pt,text=linkblue] {$d_T(x,z')\le2j+1$} (zp);
\path (0,-1.55) -- (4.5,-1.55);
\end{tikzpicture}
\end{minipage}\hfill
\begin{minipage}[t]{.48\linewidth}
\centering
{\small\textup{(b)} Controlling the extra matches\par}
\medskip
\begin{tikzpicture}[font=\small,
 point/.style={circle,fill=black,inner sep=1.6pt},
 path/.style={draw=black!75,line width=.6pt},
 bound/.style={draw=linkblue,dashed,line width=.7pt}]
\node[point,label=above:$x$] (x) at (0,0) {};
\node[point,label=above:$z'$] (zp) at (3.8,0) {};
\node[point,label=below:$v$] (v) at (3.8,-1.2) {};
\draw[path] (x) -- node[above] {$\le2j+1$} (zp);
\draw[path] (zp) -- node[right] {$\le j+1$} (v);
\draw[bound] (x) -- node[below left,text=linkblue] {$\le3j+2$} (v);
\path (0,-1.55) -- (4.5,-1.55);
\end{tikzpicture}
\end{minipage}
\caption{The two triangle inequalities in the compression step.
The pivot $z$ is replaced by its fixed label representative $z'$.
The diagrams show schematic graph paths; the labels are upper bounds
on their lengths. The logarithmic profile controls the enlargement
to radius $3j+2$.}
\label{fig:pivot-comparison}
\end{figure}


Write $z'_{j,k}=\rho_{i_0}(L_{i_0}(z_{j,k}))$ and define
\[
  A_j(x)=\sum_{y=(v,T)\in U_j}\mu_y
       \mathbf1_{\{d_T(x,z'_{j,\kappa_j(y)})\le2j+1\}}.
\]
For comparison, let
\[
  f_{\mu,j}(x)=\sum_{y=(v,T)\in Y}\mu_y
                       \mathbf1_{\{d_T(x,v)\le j\}}.
\]
Every assigned column counted by $f_{\mu,j}(x)$ is counted by $A_j(x)$,
by \eqref{eq:pivot-forward}.  The other columns counted by
$f_{\mu,j}(x)$ have total mass at most $2\theta$: at most $\theta$ lies
outside $Y_0$, and at most $\theta$ lies in $I_x\setminus U_j$.
Conversely, \eqref{eq:pivot-backward} bounds every column counted by
$A_j(x)$.  Hence
\begin{equation}
  f_{\mu,j}(x)-2\theta
  \le A_j(x)
  \le\sum_{y=(v,T)\in Y}\mu_y\mathbf1_{\{d_T(x,v)\le3j+2\}}.
  \label{eq:level-comparison}
\end{equation}

\emph{Combining the levels.}
Let $F=\sum_{j=0}^{h-1}\gamma_jA_j$.  By
\eqref{eq:profile-expansion}, $f_\mu=\sum_j\gamma_jf_{\mu,j}$.
Thus the lower bound in \eqref{eq:level-comparison} gives
$F\ge f_\mu-2\theta$.  For the upper bound, apply
\eqref{eq:profile-dilation} separately to each distance $d_T(x,v)$ and
sum with weights $\mu_{(v,T)}$.  Since the total mass is at most one and
$\log3/\log(h+1)\le\theta$, this yields the pointwise interval
\begin{equation}
  f_\mu-2\theta\le F\le f_\mu+\theta.
  \label{eq:unrounded-approximation}
\end{equation}

\emph{A finite encoding.}
The center $v$ of an assigned column no longer occurs in its indicator in
$A_j$.  Grouping columns by assigned slot and by $T$ gives exactly
\begin{equation}
  A_j(x)=\sum_{k=1}^{s}\sum_{T\in\mathcal T}
       w_{j,k,T}\mathbf1_{\{d_T(x,z'_{j,k})\le2j+1\}},
  \label{eq:grouped-approximation}
\end{equation}
where
\[
  w_{j,k,T}=
  \sum_{\substack{v\in X:\ (v,T)\in U_j\,,\ \kappa_j(v,T)=k}}
           \mu_{(v,T)}.
\]
Unused slots have all weights zero and are padded with a fixed realized
$L_{i_0}$-label.  Such a label exists because $X$ is nonempty.  Each
weight belongs to $[0,1]$, and there are at most
$M_0=s2^u$ weights at each level.  Round each weight down to a multiple
of $\delta=\theta/M_0$, writing
\[
  \widetilde w_{j,k,T}
   =\delta\left\lfloor\frac{w_{j,k,T}}{\delta}\right\rfloor.
\]
Let $\widetilde A_j$ be \eqref{eq:grouped-approximation} with these
rounded weights, and put $A=\sum_j\gamma_j\widetilde A_j$.  Since every
indicator is at most one,
\[
  0\le A_j(x)-\widetilde A_j(x)\le M_0\delta=\theta.
\]
The weights $\gamma_j$ sum to one, so $0\le F-A\le\theta$.
Together with \eqref{eq:unrounded-approximation}, this proves
\eqref{eq:positive-approximation}.

It remains to count one list that works for all $\mu$.  For fixed $i_0$,
the data determining $A$ are the $s$ pivot labels at each of the $h$
levels and the rounded weights in \eqref{eq:grouped-approximation}.
The label determines $z'_{j,k}$ through the fixed map $\rho_{i_0}$, so
no original pivot or column center must also be specified.  There are
at most $q^{hs}$ choices of labels.  Each rounded weight has at most
\[
  \left\lfloor\frac{M_0}{\theta}\right\rfloor+1
  \le1+\frac{M_0}{\theta}
\]
possible values, and there are at most $hM_0$ such weights.
Enumerate these choices for every $i_0\in[u]$, allowing also codes that
do not arise from a coefficient vector.  This gives a fixed finite list
of functions containing every approximant constructed above, with
\[
  Q\le u\,q^{hs}
        \left(1+\frac{s2^u}{\theta}\right)^{hs2^u}.
\]
Taking logarithms proves \eqref{eq:compression-count}.  The graphs and
the representative maps are shared data defining this list; they are
not additional choices for each approximant.
\end{proof}

The family $\mathcal T$ and all the graphs are fixed before applying
Proposition~\ref{prop:compression}.  The common index $i_0$ may depend on
the coefficient vector.  No separation assumption on $X$, and no choice
of a single distinguished member of $\mathcal T$, enters the argument.

\begin{corollary}[Signed compression]
\label{cor:signed-compression}
With the notation of Proposition~\ref{prop:compression}, put
\[
  f_\lambda=\sum_{y\in Y}\lambda_y g_y,
  \qquad
  B_1^Y=\left\{\lambda\in\R^Y:\sum_{y\in Y}|\lambda_y|\le1\right\},
\]
and
\[
  \mathcal F=\{f_\lambda:\lambda\in B_1^Y\}
            =\absconv\{g_y:y\in Y\}.
\]
There is a list of at most $Q^2$ functions which approximates every
member of $\mathcal F$ to uniform error $6\theta$.  Moreover, $B_1^Y$
can be partitioned into at most $Q^2$ nonempty clusters such that
\begin{equation}
  \|f_\lambda-f_{\lambda'}\|_{\infty,X}\le12\theta
  \label{eq:signed-cluster-diameter}
\end{equation}
whenever $\lambda,\lambda'$ belong to the same cluster.
\end{corollary}

\begin{proof}
Write $\lambda=\lambda^+-\lambda^-$, where
$\lambda_y^+=\max\{\lambda_y,0\}$ and
$\lambda_y^-=\max\{-\lambda_y,0\}$.  Each part is nonnegative and has
total mass at most one.  Apply Proposition~\ref{prop:compression} to
choose $A^+,A^-\in\mathcal A$ with
\[
  \|f_{\lambda^+}-A^+\|_{\infty,X}\le3\theta,
  \qquad
  \|f_{\lambda^-}-A^-\|_{\infty,X}\le3\theta.
\]
Then $A^+-A^-$ approximates $f_\lambda$ to error at most $6\theta$.
There are at most $Q^2$ differences in $\mathcal A-\mathcal A$.
Order this finite list and assign each $\lambda\in B_1^Y$ to its first
approximant with error at most $6\theta$.  The nonempty fibers form the
claimed partition, and the triangle inequality proves
\eqref{eq:signed-cluster-diameter}.  The approximating functions need
not belong to $\mathcal F$; the clusters themselves consist of actual
coefficient vectors in $B_1^Y$.
\end{proof}
